\documentclass[11pt]{article}

\usepackage[margin=1in]{geometry}
\usepackage{amsmath,amssymb}
\usepackage{booktabs}
\usepackage{array}
\usepackage{graphicx}
\usepackage{hyperref}
\usepackage{xcolor}
\usepackage{enumitem}
\usepackage{multicol}
\usepackage{caption}
\usepackage{longtable}

\hypersetup{colorlinks=true, linkcolor=blue, citecolor=blue, urlcolor=blue}

\title{\textbf{STEMS: Safe Multi-Agent Reinforcement Learning for \\
Buildings, Heat Pumps and Electric Vehicles} \\[6pt]
\large Replication, Mathematics, and Experimental Record on Real CityLearn Data}
\author{}
\date{\today}

\newcommand{\R}{\mathbb{R}}

\begin{document}
\maketitle

\begin{abstract}
This report is a complete record of a from-scratch reimplementation of STEMS
(Spatial-Temporal Enhanced Multi-Agent Safe Building Energy Management System)
\cite{stems2025} and of the extensions built on top of it. It documents, in
order: the real CityLearn environments used and every device equation they
expose; the full mathematical formulation of the architecture, from the
building-similarity graph through the spatial-temporal encoder to the
constrained actor-critic update; a unified control-barrier framework covering
the battery, the hot-water tank and electric-vehicle charging as instances of
one object; the evaluation metrics; and each experiment that was run -- including a sweep of a shared grid
constraint across a fleet of deadline-constrained vehicles -- with its
measurements, its limitations, and -- where a result was negative or a mechanism
inactive -- the specific work planned to address it. A dedicated section
describes standalone operation of the heat-pump and hot-water controller and
sets out what would be required to run it on a commercial heat-pump platform.
All measurements reported here were taken on the real simulator; anything
synthesised or derived is marked as such where it appears.
\end{abstract}

\tableofcontents
\newpage

% =====================================================================
\section{Introduction and Scope}
\label{sec:intro}
% =====================================================================

Buildings account for a large and increasingly electrified share of grid demand.
Batteries, heat pumps, hot-water stores and electric vehicles are all
controllable, all constrained, and all coupled through the same electrical
connection. Coordinating them is a constrained sequential decision problem: an
agent must reduce cost, peak demand and emissions while keeping every device
inside its physical limits and keeping occupants comfortable. Safe reinforcement
learning -- RL augmented with an explicit safety layer -- is the approach taken by
the paper this work reimplements \cite{stems2025}.

\subsection{What this report covers}

The work proceeded in four stages, and this report follows that order:

\begin{enumerate}[nosep]
\item \textbf{Replication on real data.} A faithful reimplementation of the
STEMS architecture, run on a real eight-building CityLearn simulation, with
every device parameter read from the simulator rather than assumed
(Sections~\ref{sec:citylearn}--\ref{sec:reward}).
\item \textbf{Calibration and safety experiments} on the battery-side control
barrier (Section~\ref{sec:exp-battery}).
\item \textbf{Thermal extension} -- hot water and the heat pump, including
anticipatory pre-heating and weather-dependent efficiency
(Sections~\ref{sec:exp-dhw}--\ref{sec:exp-stress}), with standalone operation and
platform integration treated separately in Section~\ref{sec:standalone}.
\item \textbf{Electric vehicles} -- unifying the hot-water and vehicle problems
under one barrier, and constructing an environment in which buildings and
vehicles coexist (Sections~\ref{sec:exp-ev}--\ref{sec:exp-merged}).
\end{enumerate}

\subsection{Reporting conventions}

Three conventions are used throughout, and are worth stating because they shape
how the numbers should be read.

\begin{itemize}[nosep]
\item \textbf{Measured, not assumed.} Every device rate, capacity and efficiency
used by the controller is read from the live simulator object. Where a quantity
is estimated rather than measured, the estimator is described and its inputs are
restricted to information available at that timestep.
\item \textbf{Fail loud.} A missing observation, an absent device or an
unsupported configuration raises an error rather than resolving silently to a
default. Simulator corrections are announced and stamped into the run metadata.
\item \textbf{Negative results are reported in place.} Where a mechanism had no
effect, or improved one quantity at the expense of another, the measurement is
given as observed and followed by the specific planned work
(Section~\ref{sec:limitations}).
\end{itemize}

% =====================================================================
\section{The CityLearn Environments}
\label{sec:citylearn}
% =====================================================================

CityLearn (2.6.0b1) is an open-source Gymnasium building-energy simulator. It
exposes per-building observations and actions and advances the simulation using
physically motivated device models driven by historical weather and load data.

\subsection{The building schema}

\begin{center}
\begin{tabular}{ll}
\toprule
Parameter & Value \\
\midrule
Source & NREL ResStock AMY2018--2021, Travis County, Texas \\
Buildings simulated & 8 (of 100 candidates), decentralised \\
Horizon & steps 0--8759 at 3600\,s $\Rightarrow$ one full year, hourly \\
Devices per building & bidirectional heat pump (cooling + heating), \\
 & electric DHW heater, battery, DHW storage tank, PV \\
Sizing & \texttt{autosize: true}; devices sized at \texttt{reset()} \\
Observations & 45 active; 28 selected by name, $+2$ in heat-pump mode \\
Actions & \texttt{[dhw\_storage, electrical\_storage, cooling\_or\_heating\_device]} \\
\bottomrule
\end{tabular}
\end{center}

Observations are mapped from CityLearn's names to the STEMS subset \emph{by
name}, so the selection is explicit and auditable rather than positional.

\subsection{Device equations}

\paragraph{Storage (battery and DHW tank).} Both inherit the same base dynamics.
For requested energy $e$ (positive = charge):
\begin{equation}
\text{SOC}^{\text{init}} = \max\!\big(0,\; \text{SOC}_{t-1}\, C\, (1-\lambda)\big),
\label{eq:soc-init}
\end{equation}
\begin{equation}
\text{SOC}_t = \frac{1}{C}\begin{cases}
\min\big(\text{SOC}^{\text{init}} + e\sqrt{\eta},\; C\big) & e \geq 0, \\[4pt]
\max\big(0,\; \text{SOC}^{\text{init}} + e/\sqrt{\eta}\,\big) & e < 0,
\end{cases}
\label{eq:soc-update}
\end{equation}
with capacity $C$, round-trip efficiency $\eta$ and standby loss $\lambda$.

\paragraph{Heat pump.} A temperature-dependent Carnot coefficient of performance
is computed for each mode and clipped to $[0,20]$:
\begin{equation}
\text{CoP}_{\text{heat}}(T_{\text{out}}) = \eta\,\frac{T^{\text{heat}}_{\text{tgt}} + 273.15}{T^{\text{heat}}_{\text{tgt}} - T_{\text{out}}},
\qquad
\text{CoP}_{\text{cool}}(T_{\text{out}}) = \eta\,\frac{T^{\text{cool}}_{\text{tgt}} + 273.15}{T_{\text{out}} - T^{\text{cool}}_{\text{tgt}}},
\label{eq:cop}
\end{equation}
\begin{equation}
P_{\text{out}} = \min\!\big(|a|\,P_{\text{nom}},\, P_{\text{nom}}\big)\cdot \text{CoP}, \qquad
P_{\text{in}} = P_{\text{out}} / \text{CoP}.
\label{eq:hp-power}
\end{equation}
Which mode is active is \emph{not} selected by the sign of the action: an
exogenous, data-derived \texttt{hvac\_mode} schedule gates it, and the scalar
action scales whichever device that schedule has enabled. This is why the
comfort term must be season-aware (Section~\ref{sec:reward}).

\paragraph{DHW heater.} A resistive model,
$P_{\text{out}} = \min(|a| P_{\text{nom}}, P_{\text{nom}})\,\eta$ with
$\eta \approx 0.9$--$0.99$, charging the tank through
Eqs.~\ref{eq:soc-init}--\ref{eq:soc-update}.

\paragraph{PV.} Generation is read from a PySAM PVWatts simulation executed once
at \texttt{reset()}, constrained by each building's roof area.

CityLearn ships a default reward; STEMS does not use it and computes its own
(Section~\ref{sec:reward}).

\subsection{An environment containing both buildings and vehicles}
\label{sec:merged-schema}

Neither shipped dataset expresses the full problem. The Travis schema has no
chargers. The EV challenge dataset has chargers but \emph{no thermal
observations at all} -- no indoor temperature, setpoints, hot water, occupancy or
cooling demand -- so it cannot represent the building side.

\texttt{setup\_citylearn\_ev.py} constructs a merged schema,
\texttt{tx\_travis\_8b\_ev}: the Travis buildings unchanged, plus charger
hardware, vehicle battery definitions and per-charger schedules.

\begin{center}
\begin{tabular}{ll}
\toprule
Buildings & 8, of which \textbf{6 have a charger} \\
Observation dimension & 35 $=$ 28 base $+$ 2 heat-pump $+$ 5 EV \\
Action dimension & 4: DHW, battery, heat pump, EV charger \\
Charger hardware & 11\,kW / 7.4\,kW, $\eta = 0.95$, 7.2\,kW discharge \\
Vehicle batteries & 40--75\,kWh \\
Charging lead time & \textbf{5.7--7.4\,h} \\
\bottomrule
\end{tabular}
\end{center}

No building data is copied. CityLearn resolves a charger file as
\texttt{os.path.join(root\_directory, charger\_simulation)}, and
\texttt{os.path.join} returns an absolute right-hand side unchanged, so the
charger schedules live in the project directory while every building file
continues to resolve from the untouched simulator cache. The thermal side of the
merged schema is therefore identical to the original.

Partial penetration (6 of 8) is deliberate: it leaves two buildings that can
only \emph{offer} flexibility and never demand it.

\paragraph{Provenance.} The building thermal models, loads, weather, pricing and
carbon intensity are real and unmodified. The \textbf{vehicle schedules are
synthetic} -- ResStock carries no vehicle data -- generated from a seeded,
documented commuter model: weekday departure $\sim$07:30 and return $\sim$17:30,
weekend 09:30/15:30, stay-at-home probability 12\%/45\%, required state of
charge at departure $\mathcal U(0.70, 0.90)$, arrival state of charge
$\mathcal U(0.25, 0.55)$ decaying with time away. Every result measured on this
schema carries that caveat.

\subsection{Heterogeneous buildings}

CityLearn buildings need not agree with one another: in the EV challenge data
the native action dimension is 1, 2 or 3 and the observation dimension 28, 35,
37 or 42, because only some buildings own a charger and charger identifiers are
embedded in observation names.

The architecture requires one shape for all buildings, since the graph network
stacks them into a single matrix. The wrapper therefore builds a \emph{canonical
padded layout} together with a per-building map into native indices:

\begin{itemize}[nosep]
\item \textbf{Observations} $=$ base features $+\;5 \times$ (EV slots). A
building with fewer bays is zero-padded, and a padded slot reports
\texttt{connected\_state} $=0$, which the vehicle barrier already interprets as
"no vehicle". An absent charger and an empty bay therefore follow the same code
path.
\item \textbf{Actions} follow a canonical order restricted to devices some
building owns. Slots a building lacks are \emph{removed} before the vector
reaches the simulator, because CityLearn sizes its action array to the devices
that exist.
\end{itemize}

On the homogeneous building schema every map is the identity, so this reduces
exactly to the original behaviour -- verified by test and by re-running a full
study (Section~\ref{sec:exp-regression}).

% =====================================================================
\section{Architecture and Mathematical Formulation}
\label{sec:stems}
% =====================================================================

Per timestep, for each building $i$:
$$
o_i \;\xrightarrow{\text{normalise}}\;
\underbrace{\text{GCN} \oplus \text{Transformer}}_{\text{shared encoder}}
\;\to\; r_i \in \R^{64} \;\xrightarrow{\pi_\theta}\; a_i^{\text{raw}}
\;\xrightarrow{\text{barrier}}\; a_i^{\text{safe}} \;\to\; \text{CityLearn}.
$$

\subsection{Building similarity graph}

Buildings are connected by a graph combining geographic proximity and functional
similarity:
\begin{equation}
w_{ij} = \alpha\, \exp\!\Big(-\frac{d_{ij}^2}{2\sigma_d^2}\Big) + \beta\, \exp\!\Big(-\frac{\|f_i - f_j\|^2}{2\sigma_f^2}\Big), \qquad w_{ii} = 0,
\label{eq:graph}
\end{equation}
with $d_{ij}$ the geographic distance and $f_i$ a normalised functional feature
vector (battery capacity, floor area).

\subsection{Spatial encoder}

A three-layer graph convolutional network propagates information across the $B$
buildings using the symmetrically normalised adjacency $\hat A = A + I$:
\begin{equation}
H^{(l+1)} = \sigma\!\Big(\hat D^{-1/2} \hat A\, \hat D^{-1/2}\, H^{(l)}\, W^{(l)}\Big), \qquad \hat D_{ii} = \sum_j \hat A_{ij},
\label{eq:gcn}
\end{equation}
with $\sigma = \mathrm{ReLU}$ and hidden width 64.

\subsection{Temporal encoder}

A multi-head self-attention block attends over a rolling $T=24$\,h window of each
building's own history:
\begin{equation}
e_{\tau,\tau'} = \frac{Q_\tau K_{\tau'}^\top}{\sqrt{d_k}}, \qquad
z_{i,t} = \mathrm{LayerNorm}\Big(x_{i,t} + \sum_{\tau'} \mathrm{softmax}_{\tau'}(e_{t,\tau'})\, V_{\tau'}\Big),
\label{eq:transformer}
\end{equation}
with 4 heads, a 32-dimensional embedding and fixed sinusoidal positional
encodings.

\subsection{Fusion and per-building actor-critic}

\begin{equation}
r_i = W_s h_i + W_t z_i + b.
\label{eq:fusion}
\end{equation}
Each building has its own squashed-Gaussian actor, value critic and multi-head
cost critic conditioned on $r_i$:
\begin{align}
z &\sim \mathcal N(\mu(r_i), \sigma(r_i)^2), \quad a = \tanh(z), \\
\log \pi_\theta(a\mid r_i) &= \log \mathcal N(z;\mu,\sigma^2) - \sum_d \log\!\big(1 - \tanh^2(z_d) + \epsilon\big), \\
V_\phi(r_i) &\in \R, \qquad V^c_\psi(r_i) \in \R^K \ (K=3).
\end{align}

\subsection{Constrained update}

Training is on-policy, one update per full-year episode. Critics are trained by
TD regression against target networks; the actor by a generalised advantage
estimate ($\gamma = 0.99$, $\lambda = 0.95$) penalised by the Lagrangian cost
advantage, with an auto-tuned entropy temperature (target entropy
$-|\mathcal A|$):
\begin{align}
\mathcal L_{\text{critic}} &= \mathbb E\big[(V_\phi(r_i) - (r + \gamma V_\phi^{\text{tgt}}(r_i')))^2\big], \\
\mathcal L_{\text{cost}} &= \mathbb E\big[(V^c_\psi(r_i) - (c + \gamma V^{c,\text{tgt}}_\psi(r_i')))^2\big], \\
\hat A_i^{\text{eff}} &= \hat A_i^{\text{GAE}} - \textstyle\sum_k \max(\lambda_k,0)\, \hat A^{c}_{i,k}, \\
\mathcal L_{\text{actor}} &= -\mathbb E\big[\hat A_i^{\text{eff}}\, \log\pi_\theta(a_i\mid r_i)\big] + \alpha\, \mathbb E[\log\pi_\theta(a_i\mid r_i)].
\end{align}

% =====================================================================
\section{Safety: One Barrier Framework, Four Devices}
\label{sec:safety}
% =====================================================================

\subsection{Classical barriers}

Three barriers are evaluated on the \emph{next predicted} state:
\begin{align}
h_1&: \text{SOC} \in [\text{SOC}_{\min}, \text{SOC}_{\max}] = [0.1, 0.9], \label{eq:h1}\\
h_2&: |e_i| \le P_{\text{building,max}}, \label{eq:h2}\\
h_3&: \textstyle\sum_i \max(e_i, 0) \le P_{\text{grid,max}}. \label{eq:h3}
\end{align}

Because $h_1$ is decoupled across buildings it is solved \emph{analytically} per
building rather than through a general quadratic program -- exact, and fast
enough for the $8760 \times 8$ step budget. With the per-building state-of-charge
change per unit action $\rho_i = P_{\text{nom},i}/C_i$ read from the real
battery, a robust margin $m$ and an anticipatory horizon $H$:
\begin{equation}
[\ell_i, u_i] = \Big[\text{SOC}_{\min} + m + \tfrac{1}{2}\rho_i H,\;\; \text{SOC}_{\max} - m - \tfrac{1}{2}\rho_i H\Big],
\label{eq:enforced-band}
\end{equation}
\begin{equation}
a^{\text{batt,safe}}_i =
\begin{cases}
\mathrm{clip}\Big(a^{\text{nom}}_i,\; \tfrac{\ell_i - \text{SOC}_i}{\rho_i},\; \tfrac{u_i - \text{SOC}_i}{\rho_i}\Big) & \text{if feasible}, \\[6pt]
+1 & \text{if } \text{SOC}_i < \ell_i \ \text{(recover: charge)}, \\[2pt]
-1 & \text{if } \text{SOC}_i > u_i \ \text{(recover: discharge)}.
\end{cases}
\label{eq:cbf-projection}
\end{equation}
The projection is feasible by construction: a recovery action exists even when
the state is already outside the band, so the shield never falls back to an
all-zero emergency action.

\subsection{Lagrangian safety pressure}

Alongside the hard shield, a soft penalty trains the \emph{policy} to be safe,
with one multiplier per constraint. With cost-rate error
$\epsilon_k = J_{c_k} - d$ and target rate $d$:
\begin{equation}
I_k \leftarrow \mathrm{clip}(I_k + k_i \epsilon_k,\, 0,\, \lambda_{\max}), \qquad
\lambda_k \leftarrow \mathrm{clip}\big(k_p \epsilon_k + I_k + k_d\, \mathrm{relu}(J_{c_k}^{(t)} - J_{c_k}^{(t-1)}),\, 0,\, \lambda_{\max}\big).
\label{eq:pid-lagrangian}
\end{equation}

\subsection{The unification: deadline-constrained storage}
\label{sec:deadline}

Three of the four controllable devices are the same mathematical object:

\begin{center}
\emph{a store that must reach a required state of charge by a deadline, while
charging at a finite rate.}
\end{center}

\begin{table}[h]
\centering
\caption{The same barrier, three devices.}
\label{tab:unification}
\small
\begin{tabular}{lccc}
\toprule
 & DHW tank & Electric vehicle & Battery \\
\midrule
Deadline & implicit: horizon $L$ & \textbf{declared} & none (a band) \\
Requirement & forecast demand $/\,C$ & \textbf{declared} & $\text{SOC}_{\min}$ \\
Capacity & fixed & \textbf{varies per arrival} & fixed \\
Rate $\rho$ & $\min(c, P\eta/C)$ & $P_{\text{ch}}\eta / C_{\text{batt}}$ & $P_{\text{nom}}/C$ \\
Measured lead time & 1.2--1.9\,h & 5.7--7.4\,h & --- \\
Slack & always 0 & \textbf{positive} & --- \\
\bottomrule
\end{tabular}
\end{table}

\texttt{DeadlineStorageBarrier} implements the shared mechanism once; the
hot-water and vehicle barriers are subclasses supplying only a requirement
function. This is a discrete-time, rate-limited instance of input-constrained
control-barrier methods, in which safety under actuator limits is certified by
predicting forward under a backup control (here, "charge at maximum") over a
finite horizon.

\paragraph{Horizon rule.} With gap $g = s^{\text{req}} - s$ and rate $\rho$,
\begin{equation}
\text{steps\_needed} = \lceil g / \rho \rceil, \qquad
\text{slack} = \text{steps\_to\_deadline} - \text{steps\_needed},
\label{eq:horizon-rule}
\end{equation}
and $\text{slack} \le 0$ is a \emph{latest-start} trigger: charge now or miss.
The hot-water tank passes $\text{steps\_to\_deadline} = 0$, so every unmet
requirement is immediately urgent; a vehicle passes its real time to departure,
which is what creates deferral room.

\paragraph{Projection.} Monotone -- it raises the action and never lowers one, so
it can add readiness but can never veto a policy that already charges harder:
\begin{equation}
a^{\text{safe}}_i = \max\Big(a^{\text{nom}}_i,\ \min\big(c_i,\ \tfrac{\max(s^{\text{req}}_i - s_i,\,0)}{\rho_i}\,c_i\big)\Big).
\label{eq:deadline-projection}
\end{equation}

\subsection{Hot water: estimating the requirement}
\label{sec:dhw-forecast}

CityLearn exposes hot-water demand for the current step only; there is no demand
forecast observation. The requirement is therefore estimated online as an
hour-of-day climatology over already-observed steps, per building $i$, hour $h$
and weekday/weekend flag $w$:
\begin{equation}
m_{i,h,w} \leftarrow m_{i,h,w} + \alpha\,(d_{i,t} - m_{i,h,w}),
\label{eq:ewma}
\end{equation}
\begin{equation}
\hat D_i(L) = \sum_{k=1}^{L} m_{i,(h+k)\bmod 24,\,w}\,\big(1 + g\,\mathrm{relu}(\bar T - \hat T_{t+k})\big),
\label{eq:dhw-forecast}
\end{equation}
with a cold-weather uplift. No future information enters: $\hat T_{t+k}$ comes
from the 1/2/3-step outdoor temperature predictions already present in the
observation vector, and $m$ is built only from past steps. During warm-up the
estimator falls back to $L \cdot d_{i,t}$, which never reports less demand than
is occurring. The barrier is
\begin{equation}
h_4:\quad \text{SOC}^{\text{dhw}}_i \;\geq\; \mathrm{clip}\!\left(\frac{\hat D_i(L)}{C_i} + m_{\text{dhw}} + \kappa\, \delta^{\text{cop}}_i,\ 0,\ s_{\max}\right).
\label{eq:h4}
\end{equation}

\subsection{Weather-dependent efficiency}
\label{sec:weather}

Weather enters twice.

\paragraph{Efficiency-aware power barriers.} Space conditioning draws
$|a| P_{\text{nom}}$ electrical, and the thermal service that buys scales with
the coefficient of performance, so the same comfort costs more power on a cold
hour. The guard counts that draw and derates the enforced cap in proportion to
the efficiency shortfall against a rated reference:
\begin{equation}
P^{\text{eff}}_{\text{max},i} = P_{\text{max}}\Big(1 - \mathrm{clip}\big(1 - \tfrac{\text{CoP}_i(T_{\text{out}})}{\text{CoP}_i(T_{\text{ref}})},\,0,\,0.5\big)\Big).
\label{eq:cop-derate}
\end{equation}

\paragraph{Cold-front anticipation.} A falling temperature forecast means both
higher upcoming demand and lower efficiency, so it is worth storing heat before
the front arrives:
\begin{equation}
\delta^{\text{cop}}_i = \mathrm{clip}\!\left(\frac{\text{CoP}_i(T_t) - \min_{k \le H}\text{CoP}_i(\hat T_{t+k})}{\text{CoP}_i(T_t)},\ 0,\ 1\right),
\label{eq:cop-drop}
\end{equation}
which is identically zero under steady or improving weather.

\subsection{Coupled feasibility}
\label{sec:coupled}

Each barrier is individually feasibility-guaranteed. That property does not
survive coupling. Under a shared power cap, meeting every deadline requires
\begin{equation}
\sum_i \frac{(s^{\text{req}}_i - s_i)\,C_i}{\eta_i} \;\le\; P_{\text{cap}}\, \Delta t\, T,
\label{eq:coupled-feasibility}
\end{equation}
with $T$ the time to the earliest binding deadline. When this fails the safe set
is \emph{empty}: no action sequence satisfies every deadline and the cap
simultaneously, and no per-device projection can repair it.

The implementation detects the condition rather than absorbing it.
\texttt{coupled\_feasibility()} evaluates Eq.~\ref{eq:coupled-feasibility};
\texttt{prioritise()} allocates the available power earliest-deadline-first, so
that when something must slip the choice is an explicit policy rather than an
artifact of array ordering; and \texttt{CBFShield.feasibility\_report()} surfaces
both. Reporting an unreachable requirement is the alternative to claiming a
guarantee that does not hold.

% =====================================================================
\section{Reward and Metrics}
\label{sec:reward}
% =====================================================================

The reward decomposes additively into four terms per building per timestep:
\begin{align}
r^{\text{econ}}_i &= -\mu\, v_t\, e_i, \label{eq:r-econ}\\
r^{\text{stab}}_i &= \alpha_g\Big(1 - \tfrac{E_{\text{grid}}}{P_{\text{grid,max}}}\Big)^2
+ \alpha_b\Big(1 - \tfrac{|e_i|}{P_{\text{building,max}}}\Big)
- \beta_r \tfrac{|e_i - e_i^{\text{prev}}|}{P_{\text{building,max}}}, \label{eq:r-stab}\\
r^{\text{comfort}}_i &= -\lambda_c \cdot d(T^{\text{in}}_i)^2, \qquad
d(T^{\text{in}}) = \begin{cases}
T^{\text{in}} - T^{\text{cool}} & T^{\text{in}} > T^{\text{cool}} \\
T^{\text{heat}} - T^{\text{in}} & T^{\text{in}} < T^{\text{heat}} \\
0 & \text{otherwise}
\end{cases}, \label{eq:r-comfort}\\
r^{\text{renew}}_i &= \xi \cdot \min\!\Big(\tfrac{\text{solar}_i}{\max(\text{load}_i,\epsilon)},\, 1\Big). \label{eq:r-renew}
\end{align}

The dual-setpoint deadband in Eq.~\ref{eq:r-comfort} is a correctness
requirement rather than a refinement: with a single cooling setpoint the agent is
penalised for being comfortably warm in winter. Any control mode that drives the
heat pump enables it automatically.

\paragraph{Metrics.} Seven quantities are computed from \emph{observed} post-step
state, never from the barrier's internal model, so the safety figure cannot be
influenced by the shield's own assumptions: cost, emissions, average daily peak,
total consumption, ramping rate, discomfort rate, and safety violation rate.
Metrics 1--5 are reported normalised against a rule-based reference controller,
following the convention of the metric definition; 6--7 are absolute. Violations
are additionally decomposed into \emph{avoidable} and \emph{unavoidable} from the
pre-action state, which distinguishes a control failure from a physically forced
one. Hot-water experiments add readiness rate, coverage rate, energy deficit and
mean tank state of charge.

% =====================================================================
\section{Experiments}
\label{sec:experiments}
% =====================================================================

\subsection{E1 --- Battery calibration and the safety ablation}
\label{sec:exp-battery}

The dominant factor in the measured safety-violation rate is how precisely the
barrier's state-of-charge model matches each building's actual battery. Reading
the real per-building dynamics gives
$\rho_i = P_{\text{nom},i}/C_i \in [0.178, 0.533]$ -- a threefold variation across
eight buildings. An earlier implementation used a single constant $0.1$ for
every building.

The same seeded nominal policy was replayed through the real environment under
different safety configurations, which isolates the shield's marginal effect
from training noise.

\begin{table}[h]
\centering
\caption{E1: safety configuration ablation (800 steps $\times$ 2 seeds).}
\label{tab:ablation}
\begin{tabular}{lcc}
\toprule
Safety configuration & Violation rate & Cost \\
\midrule
No barrier (raw policy) & 0.7195 & 7676 \\
Barrier, single uniform rate estimate & 0.0945 & 7777 \\
Barrier $+$ real per-building calibration & 0.0110 & 7703 \\
\quad $+$ robust margins & 0.0000 & 7705 \\
\quad $+$ anticipatory buffer & 0.0000 & 7662 \\
\textbf{Full configuration} & \textbf{0.0000} & \textbf{7649} \\
\bottomrule
\end{tabular}
\end{table}

Three observations. Calibration precision accounts for most of the improvement:
a uniform rate estimate leaves a 9.45\% violation rate even with the shield
active, because it under-projects the safe action relative to what each battery
actually does. Per-building calibration alone reaches $0.0110$. The full
configuration reaches zero measured violations and simultaneously the lowest
cost, because on this dataset the power and grid limits never bind -- loads are
10--40\,kW against far larger caps -- so the state-of-charge barrier is the only
active constraint and it is now solved exactly.

\subsection{E2 --- Converged full-year training run}
\label{sec:exp-converged}

A 15-episode run over full-year (8760-step) episodes, seed 0, with the thermal
actuators under control.

\begin{center}
\begin{tabular}{lcc}
\toprule
Metric & Reference controller & STEMS \\
\midrule
cost & 34008.88 & 15219.87 \\
emissions & 21071.71 & 8688.99 \\
average daily peak & 36.78 & 16.47 \\
electricity consumption & 142816.61 & 55219.75 \\
ramping rate & 6.368 & 4.804 \\
discomfort rate & 0.0001 & 0.0001 \\
safety violation rate & 0.6465 & \textbf{0.0000} \\
\bottomrule
\end{tabular}
\end{center}

Episode reward rose from 205 to 671 between episodes 2 and 15 while the measured
violation rate remained at $0.0000$ in \emph{every} episode: the safety property
held throughout learning, not only at convergence.

\paragraph{Limitation.} This run predates the hot-water correction described in
Section~\ref{sec:defects}, so it is a valid \emph{heat-pump-only} result rather
than a joint hot-water and heat-pump one. The planned remedy is
Section~\ref{sec:limitations}, item~2.

\subsection{E3 --- Hot-water lead time}
\label{sec:exp-leadtime}

CityLearn charges the tank by $e = a\,C$ but clamps that to the heater's
one-hour output, so the achievable state-of-charge gain in one step is
\begin{equation}
\rho^{\text{dhw}}_i = \min\!\Big(c_i,\ \frac{P_{\text{nom},i}\,\eta_i}{C_i}\Big),
\qquad \tau_i = 1/\rho^{\text{dhw}}_i \ \text{hours}.
\label{eq:dhw-rate}
\end{equation}

Measured across the eight buildings: $\rho^{\text{dhw}} \in [0.52, 0.85]$ and
$\tau \in [1.17, 1.92]$\,h, mean $1.60$\,h. Equation~\ref{eq:dhw-rate} was
validated against the simulator rather than assumed: charging an empty tank at
maximum action for one step moves the state of charge to within 15\% of
$\rho^{\text{dhw}}_i$ for every building, the residual being tank efficiency and
standby loss.

Since $\tau_i > 1$\,h everywhere, a controller that responds to demand
\emph{after} observing it is structurally too late. This is the measurement that
motivates a pre-heat horizon $L \ge 2$.

\subsection{E4 --- Anticipatory pre-heating}
\label{sec:exp-dhw}

1500 steps $\times$ 2 seeds over the heating season. The nominal controller is
\emph{price-aware but reactive}: it refills the tank only once demand has
appeared and declines to heat above its running-mean price. Measuring against a
controller that is already trying to be economical makes the barrier's
contribution interpretable as the value of anticipation specifically. Every
configuration is graded by the \emph{same} fixed evaluation barrier, a separate
instance from the one being varied.

\begin{table}[h]
\centering
\caption{E4: anticipatory pre-heating. Readiness and coverage higher is better;
deficit, cost and \$/kWh lower is better.}
\label{tab:preheat}
\small
\begin{tabular}{lccccc}
\toprule
Configuration & Readiness & Covered & Deficit [kWh] & Cost & DHW \$/kWh \\
\midrule
No pre-heat (reactive) & 0.907 & 0.939 & 1026.5 & 3021.6 & 0.2200 \\
\quad $+$ efficiency-aware power guard & 0.907 & 0.939 & 1026.5 & 3021.6 & 0.2200 \\
Pre-heat $L=1$ & 0.928 & 0.950 & 661.7 & 3025.9 & 0.2300 \\
Pre-heat $L=2$ & 0.961 & 0.958 & 90.6 & 3033.1 & 0.2411 \\
Pre-heat $L=3$ & 0.989 & 0.968 & 33.8 & 3039.0 & 0.2936 \\
\quad $+$ cold-front anticipation ($L=2$) & 0.982 & 0.964 & 67.3 & 3035.6 & 0.2517 \\
\textbf{Full ($L=2$, weather $+$ margin)} & \textbf{0.991} & \textbf{0.970} & \textbf{47.4} & 3038.5 & 0.2571 \\
\bottomrule
\end{tabular}
\end{table}

\paragraph{Positive result 1: the horizon is set by the physics.} Readiness rises
monotonically with $L$ and the energy deficit falls by 86\% between $L=1$ and
$L=2$ ($661.7 \to 90.6$\,kWh), improving only slowly thereafter. This is the
discontinuity Eq.~\ref{eq:dhw-rate} predicts: with a mean time-to-heat of
$1.60$\,h a one-hour horizon cannot in general close the gap in time, and a
two-hour horizon can. The horizon is determined by measured device physics, not
selected by preference -- the same lesson as the battery calibration in E1,
applied to a different actuator.

\paragraph{Positive result 2: targeted anticipation is cheaper than uniform
conservatism.} Comparing $L=3$ against the full configuration ($L=2$ plus the
cold-front term and a wider margin): the full configuration reaches higher
readiness ($0.991$ vs $0.989$) and better coverage while paying 12\% less per
kWh of hot-water heating ($0.2571$ vs $0.2936$). Extending the horizon raises the
requirement at \emph{every} hour; conditioning on the forecast raises it only
when the weather signal justifies it, since Eq.~\ref{eq:cop-drop} is identically
zero under steady weather.

\paragraph{Negative result 1: readiness is not free.} Every configuration that
improves readiness also spends more: hot-water electricity rises from $196.2$ to
$238.4$\,kWh and the mean price paid per kWh from $0.2200$ to $0.2571$, lifting
total cost by 0.6\%. The mechanism differs from E1, where the barrier only ever
\emph{restricted} an action the policy already wanted: here the barrier
\emph{commands} heating at hours the policy would have skipped, and because the
requirement is driven by demand rather than price those hours are on average
more expensive than the ones a price-following controller selects.
\textbf{Planned work:} the barrier currently heats as late as the constraint
allows. Selecting \emph{when within the horizon} to pre-heat -- the earliest
inexpensive hour rather than the last feasible one -- should recover most of the
premium and turn the barrier into a load-shifting mechanism
(Section~\ref{sec:limitations}, item~3).

\paragraph{Negative result 2: the efficiency-aware power guard is inactive.} The
first two rows are identical to every reported digit. This is expected rather
than faulty: per-building loads are 10--40\,kW against an 80\,kW cap, so $h_2$
and $h_3$ never bind, and a guard that reserves headroom inside a cap that is
never approached cannot alter any action. The row is reported rather than
omitted so that an inactive mechanism is visible as inactive.
\textbf{Planned work:} exercised under a tightened cap in E5, and to be evaluated
on a heating-dominated climate where the efficiency term is larger
(Section~\ref{sec:limitations}, item~4).

\paragraph{Note on the seed spread.} The per-seed standard deviation of readiness
is exactly $0.000$ throughout. That reflects determinism, not stability: the
reactive nominal controller is a deterministic function of the observation and
the scenario seed is fixed in the schema, so both seeds trace the same
trajectory. These seeds establish reproducibility and say nothing about
robustness. \textbf{Planned work:} Section~\ref{sec:limitations}, item~5.

\subsection{E5 --- Cold-snap stress}
\label{sec:exp-stress}

Two corrections to the experimental design were required before the scenario
tested what it was intended to test, and both are recorded because each is easy
to get wrong.

\begin{enumerate}[nosep]
\item \textbf{A constant temperature offset cannot exercise a cold-front term.}
Equation~\ref{eq:cop-drop} keys on the \emph{difference} between the current and
forecast efficiency. Shifting the current reading and all of its forecasts by the
same amount changes the temperature level but leaves that difference essentially
untouched. The correction is a forecast-only ramp,
$\hat T_{t+k} \mathrel{+}= \gamma k$ with $\gamma < 0$.
\item \textbf{The cap must actually bind.} With the battery frozen, mean
per-building draw is $0.88$\,kW; a 12\,kW cap is no more binding than the nominal
one. The cap had to be reduced to 2\,kW before $h_2$ engaged at all.
\end{enumerate}

Both perturbations act on the observation stream rather than on the simulator's
internal physics, so this scenario probes whether the anticipation logic responds
correctly to a weather signal; it is not a statement about true energy use in a
cold snap.

\begin{table}[h]
\centering
\caption{E5: cold-snap stress ($-18^\circ$C observed, $-4^\circ$C/h forecast
front, 2\,kW per-building cap; 1500 steps, seed 0).}
\label{tab:stress}
\small
\begin{tabular}{lccccc}
\toprule
Configuration & Readiness & Deficit [kWh] & Cost & DHW \$/kWh & Violation \\
\midrule
No pre-heat (reactive) & 0.895 & 1163.5 & 3021.6 & 0.2200 & 0.2402 \\
\quad $+$ efficiency-aware guard & 0.895 & 1163.5 & \textbf{3008.2} & 0.2200 & 0.2388 \\
Pre-heat $L=1$ & 0.919 & 734.7 & 3013.1 & 0.2307 & 0.2390 \\
Pre-heat $L=2$ & 0.961 & 103.2 & 3021.4 & 0.2459 & 0.2390 \\
Pre-heat $L=3$ & 0.991 & 33.2 & 3028.9 & 0.3127 & 0.2402 \\
\quad $+$ cold-front anticipation ($L=2$) & 0.986 & 60.2 & 3024.9 & 0.2550 & 0.2390 \\
\textbf{Full ($L=2$, weather $+$ margin)} & \textbf{0.991} & \textbf{43.5} & 3028.3 & 0.2591 & 0.2390 \\
\bottomrule
\end{tabular}
\end{table}

\paragraph{Positive result: anticipation is worth more when there is weather to
anticipate.} Adding the cold-front term at $L=2$ raises readiness from $0.961$ to
$0.986$ and cuts the deficit by 42\% ($103.2 \to 60.2$\,kWh), a larger gain than
the same term produced under steady weather. The full configuration matches the
readiness of $L=3$ while paying 17\% less per kWh ($0.2591$ vs $0.3127$), against
12\% under nominal conditions.

\paragraph{Positive result: the efficiency-aware guard becomes active.} With the
cap binding, the second row separates from the first: cost falls 0.44\%
($3021.6 \to 3008.2$), making it the least expensive configuration in the table,
because the guard trims conditioning draw hardest exactly when efficiency is
lowest.

\paragraph{Negative result: the guard cannot repair the constraint.} Its effect
on the violation rate is negligible ($0.2402 \to 0.2388$), and the reason
quantifies cleanly. With the battery frozen, the shield commands only the
hot-water heater and the heat pump, together $366.5$ of $10597.3$\,kWh --
\textbf{3.5\% of total consumption}. The remaining 96.5\% is non-shiftable load,
and no thermal action can bring a building under a 2\,kW cap that its baseline
load already exceeds; roughly 24\% of (timestep, building) pairs violate $h_2$ in
every configuration because those violations are unavoidable in the sense of
Section~\ref{sec:reward}. A shield can only be as effective as the fraction of
load it actually commands. \textbf{Planned work:} re-enable the battery so the
controlled fraction is large enough for the guard to act on, and evaluate the
same scenario with vehicles present, where the controllable fraction is far
larger (Section~\ref{sec:limitations}, item~1).

\subsection{E6 --- Vehicle barrier on real charger data}
\label{sec:exp-ev}

The vehicle barrier was validated against the shipped CityLearn EV dataset
before being used, in the same way the tank calibration was checked against the
real tank.

Measurements: a charger rated 11\,kW at $\eta = 0.95$ serving a 60\,kWh battery
gives $\rho = 0.174$\,h$^{-1}$, i.e.\ approximately 5.7\,h to fill -- roughly
three and a half times the hot-water lead time of E3. The test suite asserts
that every required observation is present; that $\rho$ matches
$P_{\text{ch}}\eta/C$ read from the live charger object; that a 72-step rollout
spanning genuine arrivals \emph{and} empty bays leaves the barrier finite,
in-bounds and inert on empty bays; and that an unreachable departure requirement
is reported rather than absorbed.

\paragraph{Consequence for the horizon rule.} For hot water the horizon is a
hyper-parameter validated against a measured time constant. For a vehicle the
horizon is supplied by the data itself, so Eq.~\ref{eq:horizon-rule} is exercised
with nothing to tune -- and the longer lead time makes it correspondingly more
consequential.

\subsection{E7 --- The merged building-and-vehicle environment}
\label{sec:exp-merged}

The merged schema of Section~\ref{sec:merged-schema} was validated end to end.
Result: 8 buildings, 6 with chargers, observation dimension 35, action dimension
4. All control modes resolve, and a single rollout driving the hot-water heater
and the charger simultaneously moves both stores, confirming that the thermal and
vehicle sides are live in the same environment. Charging lead times range from
5.7 to 7.4\,h.

The generator asserts that no thermal observation was lost in the merge and that
the building data files are referenced, not copied.

\subsection{E8 --- Regression verification}
\label{sec:exp-regression}

Because the hot-water and heat-pump results are used as a standalone baseline
elsewhere, every structural change to the code was verified not to move them.
Two levels of check are applied:

\begin{itemize}[nosep]
\item \textbf{Operation level.} The pre-refactor hot-water projection is retained
as a frozen reference function and differenced against the unified barrier over
200 randomised states $\times$ 3 horizons $\times$ 2 weather settings.
\item \textbf{End to end.} The complete E4 study is re-run and differenced
against the stored results after each structural change.
\end{itemize}

Across three successive refactors -- unifying the barrier, widening the
environment for heterogeneous buildings, and adding the merged schema -- the
maximum relative difference is $7.359 \times 10^{-9}$, on the average-daily-peak
metric, which is the most summation-order-sensitive quantity reported. That
residual is a one-unit-in-the-last-place \texttt{float32} effect arising from the
order of summation, and is an order of magnitude below the tolerance already
present in the readiness comparison. The corresponding test asserts agreement to
\texttt{float32} precision, not exact equality.

\subsection{E9 --- Coupling sweep with a deadline-locked vehicle fleet}
\label{sec:exp-coupling}

Sections~\ref{sec:coupled} argued that per-device feasibility does not survive a
shared connection. This experiment measures that directly on the merged
environment, and measures what a coordination rule is worth as a function of how
tightly the shared constraint binds.

\paragraph{Setting.} Six vehicles, 55.2\,kW of aggregate charger capacity, on the
eight-building merged schema. Measured beforehand: the non-vehicle building load
alone peaks at $39.6$\,kW; with every vehicle charging on demand the site peaks
at $71.5$\,kW, and a mean of 4.2 vehicles are connected simultaneously. The
shared import cap is swept from 80\,kW (slack) to 15\,kW (severely binding).

\paragraph{Arms.} The nominal controller is deliberately uncoordinated: each
household charges whatever is plugged in, which is what a distribution
transformer sees when a street electrifies. Three shield configurations are
compared.

\begin{itemize}[nosep]
\item \texttt{independent} -- each deadline barrier projects on its own. Correct
per device, collectively blind to the connection.
\item \texttt{proportional} -- the shared cap is enforced; every request is scaled
by the same factor.
\item \texttt{edf} -- the shared cap is enforced; available power is allocated
earliest-deadline-first.
\end{itemize}

The last two enforce an \emph{identical} total, so any difference between them is
attributable to the priority rule alone. All three necessarily coincide wherever
the cap is slack, which is what makes a sweep the correct instrument rather than
a single operating point.

\paragraph{Service metric.} A \emph{missed departure} is a vehicle that
disconnects with a state of charge below what it required. It is the honest
service metric here because, unlike a state-of-charge excursion, it cannot be
repaired by any later action.

\begin{table}[h]
\centering
\caption{E9: grid-cap sweep, 1500 steps, 298 departures. \texttt{miss} is missed
departures; \texttt{viol} is the fraction of steps exceeding the cap.}
\label{tab:coupling}
\small
\begin{tabular}{rccccccc}
\toprule
& \multicolumn{2}{c}{\texttt{independent}} & \multicolumn{2}{c}{\texttt{proportional}} & \multicolumn{2}{c}{\texttt{edf}} & EDF gain \\
\cmidrule(lr){2-3}\cmidrule(lr){4-5}\cmidrule(lr){6-7}
Cap [kW] & miss & viol & miss & viol & miss & viol & vs.\ proportional \\
\midrule
80 & 3 & 0.000 & 3 & 0.000 & 4 & 0.000 & $-33\%$ \\
60 & 3 & 0.013 & 1 & 0.000 & 1 & 0.000 & $0\%$ \\
45 & 3 & 0.065 & 6 & 0.001 & 5 & 0.000 & $+17\%$ \\
35 & 3 & 0.113 & 20 & 0.009 & \textbf{11} & 0.011 & $\mathbf{+45\%}$ \\
30 & 3 & 0.142 & 58 & 0.018 & \textbf{34} & 0.023 & $\mathbf{+41\%}$ \\
25 & 3 & 0.179 & 115 & 0.045 & \textbf{66} & 0.047 & $\mathbf{+43\%}$ \\
20 & 3 & 0.219 & 170 & 0.099 & \textbf{115} & 0.095 & $\mathbf{+32\%}$ \\
15 & 3 & 0.273 & 214 & 0.201 & 179 & 0.177 & $+16\%$ \\
\bottomrule
\end{tabular}
\end{table}

\begin{figure}[h]
\centering
\includegraphics[width=\textwidth]{plots/coupling_sweep.png}
\caption{E9. Left: service failure. Centre: constraint breach. Right: the
missed departures avoided by earliest-deadline-first relative to proportional
sharing, at an identical enforced total. The dotted line marks the 39.6\,kW
non-vehicle baseline peak, below which no vehicle action can bring the site
under the cap.}
\label{fig:coupling}
\end{figure}

\paragraph{Positive result 1: enforcement converts a breach into a scheduling
problem.} The \texttt{independent} arm never limits anything: it misses only 3 of
298 departures at every cap, but exceeds the cap in up to 27\% of steps and peaks
at $73.4$\,kW regardless of what the cap is set to. Enforcing the shared cap
removes almost all of that breach -- at 35\,kW the violation rate falls from
$0.113$ to $0.009$ -- and converts it into a scheduling problem, at the cost of
missed departures. This trade is the substance of the coupled-feasibility
condition of Eq.~\ref{eq:coupled-feasibility}: when the requirement set cannot be
met, something must give, and the only choice available is \emph{what}.

\paragraph{Positive result 2: the value of the priority rule is non-monotone in
coupling tightness.} This is the main finding. Comparing the two arms that
enforce the same total:

\begin{itemize}[nosep]
\item At 80 and 60\,kW the cap is slack and the rule is worth nothing.
\item Between 35 and 20\,kW it avoids 32--45\% of the missed departures that
proportional sharing incurs, peaking at $+45\%$ (11 misses against 20) at
35\,kW.
\item At 15\,kW the advantage falls back to $+16\%$.
\end{itemize}

The decline at the tight end is as informative as the peak. Once the cap is far
below what the fleet requires, most deadlines are unreachable under any
allocation, and a priority rule can only reorder failures rather than prevent
them. Coordination is therefore most valuable in an intermediate regime: it is
worthless where the constraint does not bind, and its leverage shrinks again
where the problem is infeasible for everyone. Reporting a single operating point
would have supported almost any conclusion one wished to draw.

\paragraph{Negative result: a floor set by the uncontrollable load.} Every
enforcing arm peaks at exactly $39.6$\,kW for all caps at or below 30\,kW, and
the residual violation rate rises regardless of the allocation rule. That figure
is the non-vehicle building load, which the shield does not command. Below it, no
vehicle action can bring the site under the cap, and the violations are
unavoidable in the sense of Section~\ref{sec:reward}. This is the same bound
identified in E5 from a different direction, and it quantifies what any charging
coordination scheme can and cannot deliver. \textbf{Planned work:} repeat with
the battery and thermal actuators also under control, which raises the
controllable fraction substantially and should lower that floor.

\paragraph{Negative result: the headroom estimate is conservative.} At 80\,kW the
three arms should be identical, because the realised peak of $73.4$\,kW never
reaches the cap; instead \texttt{edf} records 4 missed departures against 3. The
cause is that available headroom is computed from the previously observed net
import, which already contains that step's charging, so the allocator
occasionally engages when the cap is in fact slack. The effect is small and
conservative in the safe direction, but it is a modelling artifact rather than a
property of the system. \textbf{Planned work:} estimate the next-step
non-controllable load separately from the controllable draw, instead of using the
previous total as a proxy.

\paragraph{Scope.} This experiment characterises the coupling itself and compares
\emph{allocation rules inside the safety layer}. It is not a comparison of
learned coordination architectures: no policy is trained here, and the nominal
controller is fixed by construction so that the shield's contribution can be
isolated. Establishing where the shared constraint binds, and by how much, is a
prerequisite for that comparison -- a learned coordinator evaluated at 80\,kW
would have nothing to coordinate about -- but it is not a substitute for it. The
learned comparison is Section~\ref{sec:limitations}, item~1.

% =====================================================================
\section{Standalone Heat-Pump Operation and Platform Integration}
\label{sec:standalone}
% =====================================================================

The hot-water and heat-pump controller is designed to run without the rest of the
system. This section documents that mode and sets out what integrating it with a
commercial heat-pump platform would require.

\subsection{Running standalone}

Control is restricted by a named registry mapping a mode to a set of actuators,
resolved by device name rather than by position:
$$
\texttt{ACTION\_GROUPS} \supseteq \Big\{\;\texttt{heatpump}\!:\![\text{heat pump}],\;\;
\texttt{dhw}\!:\![\text{DHW}],\;\; \texttt{thermal}\!:\![\text{DHW},\ \text{heat pump}] \Big\}.
$$
Selecting a mode has two effects:
\begin{enumerate}[nosep]
\item \textbf{Actuator masking.} Devices outside the selected set are held at
their neutral value. For a storage device an action of zero is exactly "neither
charge nor discharge" (Eq.~\ref{eq:soc-update} with $e = 0$ leaves the state
governed only by standby loss), so freezing them is physically meaningful rather
than an arbitrary placeholder.
\item \textbf{Barrier adaptation.} The battery state-of-charge barrier is
constructed only when the battery is actually controllable, so it is skipped
rather than constraining an actuator the controller cannot move. The power
barriers remain active in every mode, since hot-water and conditioning draw still
affect the building's import.
\end{enumerate}
Requesting a device the installation does not have raises an error rather than
silently controlling nothing.

In standalone thermal mode the controller therefore consists of: the
observation encoder, the per-building policy, the hot-water readiness barrier
$h_4$, the efficiency-aware power guard, and the comfort term with its
dual-setpoint deadband. No battery, vehicle or neighbourhood coupling is
required for it to run.

\subsection{What the controller needs from a heat-pump installation}

\begin{table}[h]
\centering
\caption{Signals required for standalone thermal control.}
\label{tab:signals}
\small
\begin{tabular}{p{4.6cm}p{3.0cm}p{6.4cm}}
\toprule
Quantity & Cadence & Use \\
\midrule
Outdoor temperature, plus a 1--3\,h forecast & hourly & CoP model
(Eq.~\ref{eq:cop}); cold-front term (Eq.~\ref{eq:cop-drop}) \\
Indoor temperature and setpoints (heating and cooling) & hourly & comfort
deadband (Eq.~\ref{eq:r-comfort}) \\
Hot-water tank temperature or state of charge & hourly & barrier state in
Eq.~\ref{eq:h4} \\
Hot-water draw (measured or inferred) & hourly & requirement estimator
(Eq.~\ref{eq:ewma}); only past values are used \\
Electrical import & hourly & power barriers $h_2, h_3$ \\
Electricity price, ideally day-ahead & hourly & economic reward term \\
\midrule
Compressor rated power, tank capacity, heater rating and efficiency & once, at
commissioning & calibration of $\rho^{\text{dhw}}$ (Eq.~\ref{eq:dhw-rate}) and of
the CoP curve \\
\bottomrule
\end{tabular}
\end{table}

The control output is a scalar modulation command per actuator in $[-1, 1]$,
which maps naturally onto an inverter-driven compressor and onto a hot-water
charge request. On-off equipment would require a duty-cycle or hysteresis layer
between the barrier output and the device, which is not currently implemented.

\subsection{Integration path for a commercial platform}

Taking a Thermia installation as the concrete case -- ground- or air-source units
with an inverter-modulating compressor, simultaneous space heating and hot-water
production, and a remote monitoring service -- integration would proceed in four
steps.

\paragraph{Step 1: commissioning-time calibration.} The single most consequential
step, and the lesson of both E1 and E3. The controller must read the installed
unit's rated compressor power, tank capacity, heater rating and efficiency, and
compute $\rho^{\text{dhw}}$ and the CoP curve from them. Assuming nominal values
across an installed base is precisely the error that produced the highest
violation rates in E1. This is a one-time read per installation.

\paragraph{Step 2: the barrier runs at the edge.} The hot-water barrier
(Eq.~\ref{eq:h4}) and the power guard (Eq.~\ref{eq:cop-derate}) are analytic and
per-building; they involve no optimisation solver and no neighbourhood state.
They can run on the unit controller or a local gateway, at hourly cadence, with
negligible computation. This matters for deployment: the safety layer does not
depend on connectivity, and it never returns an undefined action -- the projection
is feasible by construction, and where a requirement is genuinely unreachable it
is reported rather than silently approximated.

\paragraph{Step 3: the demand estimator learns per installation.} The requirement
estimator (Eq.~\ref{eq:ewma}) is an online hour-of-day climatology that needs no
prior training and no data leaving the property; it converges within a few days
of observed operation, and during warm-up it falls back to a conservative
estimate that never under-reports current demand.

\paragraph{Step 4: the policy is optional.} The barrier improves hot-water
readiness on top of \emph{any} nominal controller -- E4 measured its effect on a
conventional price-aware reactive controller, not on a learned policy. An
installation can therefore adopt the safety and anticipation layer alone, with the
learned policy added later. This staging substantially reduces deployment risk.

\subsection{Recalibration required for a different market}

Three items must be revisited before any claim is made about a heating-dominated
installed base.

\begin{enumerate}[nosep]
\item \textbf{Climate.} All measurements in this report come from a
cooling-dominated Texas dataset. The efficiency mechanisms of
Section~\ref{sec:weather} were inactive at nominal caps precisely because the
climate rarely stresses them (E4, negative result 2). Heating-dominated data is
required, and is the planned next step
(Section~\ref{sec:limitations}, item~4).
\item \textbf{Ground-source thermal behaviour.} The CoP model of
Eq.~\ref{eq:cop} is parameterised by a supply and a source temperature. For a
ground-source unit the source is the brine circuit, which is far more stable
hour-to-hour than outdoor air but drifts seasonally, and over long horizons is a
shared resource between neighbouring installations. Representing that would be a
change of source-temperature model, not of the barrier.
\item \textbf{Periodic disinfection.} A thermal disinfection cycle is a hard,
periodic, regulated requirement that conflicts directly with efficiency, since
the elevated tank temperature it needs is expensive to reach. It maps onto the
deadline barrier of Section~\ref{sec:deadline} without modification: a required
state that must be met within a bounded window, at a finite heating rate. This is
a natural extension of the existing mechanism and is listed in
Section~\ref{sec:limitations}, item~7.
\end{enumerate}

% =====================================================================
\section{Defects Found and Corrected}
\label{sec:defects}
% =====================================================================

Four defects were found during this work. All four are now covered by regression
tests. They are recorded here because each one silently affected measurements
until it was found.

\paragraph{D1 --- The hot-water actuator was inert.} CityLearn's storage update
for domestic hot water scales the requested energy by the \emph{heating} storage
tank's capacity rather than the hot-water tank's. The building schema used here
has no heating storage tank, so that capacity is zero and every hot-water action
resolved to zero energy. Verified directly: commanding a charge action of $0.8$
for eight consecutive steps left the tank state of charge at exactly $0.0$. The
correction is applied per building, only where the defect is present, announced
by a banner and recorded in the run metadata, with an explicit opt-out. All
device-level hot-water readings taken before this correction are invalid; the
affected runs are identified in Section~\ref{sec:exp-converged}.

\paragraph{D2 --- Vehicle departure time misread as a clock hour.} The charger
schedule field is documented as a countdown in timesteps and the shipped data
files decrement it towards zero. An initial implementation of the vehicle barrier
treated it as an absolute hour and performed clock arithmetic on it, which would
have produced incorrect deadlines at every bay. The field is now used directly.
The associated required-state field is stored as a percentage and normalised on
load.

\paragraph{D3 --- Departure countdown truncated one step early.} The reference
data keeps a vehicle parked while the countdown reaches zero, so the departure
hour itself is a final charging opportunity. The schedule generator initially
stopped one step earlier, silently removing it.

\paragraph{D4 --- State-of-charge metric scored on a frozen battery.} Under
thermal-only control the reported state-of-charge violation rate was pinned at
$1.0$, because the metric scored the frozen, uncontrolled battery against its
band. The metric is now conditioned on whether the battery is actually
controllable; the affected mode moved from $1.0$ to $0.0$.

% =====================================================================
\section{Limitations and Planned Work}
\label{sec:limitations}
% =====================================================================

Each item states an observed limitation and the specific work intended to address
it. The ordering reflects how much each would strengthen the conclusions.

\begin{enumerate}
\item \textbf{Learned coordination has not yet been compared.} E9 measured the
shared constraint and compared allocation rules \emph{inside} the safety layer,
establishing that the useful operating regime is roughly 20--45\,kW on this
fleet. It did not train anything. \emph{Planned:} with that regime now known,
compare centralised-critic and fully decentralised training across the same
sweep, which is the comparison E9 was built to make possible. Evaluating such a
comparison at a slack cap would measure nothing, which is why the sweep came
first.

\item \textbf{The controllable fraction of load bounds every result.} E5 measured
the shield commanding 3.5\% of consumption under thermal-only control; E9 found
a hard floor at the 39.6\,kW non-vehicle baseline. \emph{Planned:} run both with
the battery and the thermal actuators under control simultaneously, which is the
only way to raise that fraction.

\item \textbf{Thermal results predate the hot-water correction.} The converged
run of E2 and the earlier control-mode validation drove an inert hot-water
actuator (defect D1), so they stand as heat-pump-only results. \emph{Planned:}
re-run both to convergence on the corrected environment to obtain a genuine joint
hot-water and heat-pump result.

\item \textbf{Readiness is bought at a cost premium.} E4 measured a 0.6\%
increase in total cost and a rise from $0.2200$ to $0.2571$ in the price paid per
kWh of hot-water heating, because the barrier heats as late as the constraint
allows and is driven by demand rather than price. \emph{Planned:} choose the
heating instant \emph{within} the available horizon rather than at its end,
selecting the least expensive feasible hour. Vehicles are the natural first
application, because their deadlines leave positive slack whereas hot water does
not.

\item \textbf{Only one climate has been evaluated.} Every measurement here is
from a cooling-dominated dataset, which is why the efficiency-aware mechanisms
were inactive at nominal caps. \emph{Planned:} repeat E4 and E5 on the
heating-dominated datasets available in the simulator, where the coefficient of
performance varies over a much wider range and the cold-front term should carry
more information.

\item \textbf{Reported spreads are not uncertainty estimates.} The per-seed
standard deviation is exactly zero because the nominal controller is
deterministic and the scenario seed is fixed. \emph{Planned:} vary the stochastic
policy seed, the weather year and the building subset, and report confidence
intervals and paired comparisons instead of raw spreads.

\item \textbf{Comparative evaluation is incomplete.} The evaluation path
currently scores against a single rule-based reference, while several other
controllers are implemented but unused. \emph{Planned:} calibrate every one of
them from the live device parameters before use -- one of them presently carries
the same uniform rate assumption identified in E1 -- and report the full metric
table. Calibrating the comparators is a precondition, not an optimisation.

\item \textbf{Periodic disinfection is not yet modelled.} A regulated thermal
disinfection cycle is a deadline constraint of exactly the form already
implemented. \emph{Planned:} add it as a further instance of the deadline
barrier, with the requirement window as the deadline and the elevated tank
temperature as the required state.

\item \textbf{Vehicle-to-grid capability is unused.} The merged schema is
already discharge-capable at 7.2\,kW. \emph{Planned:} enable bidirectional
operation, which converts a vehicle from a purely constraining load into a
flexibility source and changes the structure of the coupled feasibility condition
in Eq.~\ref{eq:coupled-feasibility}.

\item \textbf{Engineering hygiene.} The project lacks packaging, continuous
integration and a container definition, and the requirement estimator's per-step
loop dominates study runtime. \emph{Planned:} package the project, add automated
test execution, and vectorise the estimator.
\end{enumerate}

% =====================================================================
\section{Summary}
\label{sec:summary}
% =====================================================================

This report records a reimplementation of the STEMS architecture on real
building data and the extensions built on it. The recurring technical theme is
that \emph{actuator dynamics must be measured from the plant rather than
assumed}, and that anticipatory safety horizons must be derived from those
measurements: a uniform state-of-charge rate left a 9.45\% violation rate with
the shield active (E1), while per-building calibration and the associated margins
reduced the measured rate to zero at no cost penalty; the same principle applied
to hot water produced an 86\% reduction in energy deficit at exactly the horizon
the measured 1.60\,h time-to-heat predicts (E4).

The safety layer was then generalised: the hot-water tank, the vehicle battery
and the electrical battery are one object -- a store that must reach a required
state by a deadline at a finite rate -- and are implemented once
(Section~\ref{sec:deadline}). That generalisation also exposed the limit of the
guarantee: individual barriers are feasible by construction, but coupled
deadlines under a shared cap can render the safe set empty
(Eq.~\ref{eq:coupled-feasibility}), which the implementation detects and reports
rather than absorbing.

The thermal controller runs standalone and is deployable independently of the
rest of the system: its barrier is analytic, per-building, solver-free and
requires no neighbourhood state, and its demand estimator learns on-site from
observed operation (Section~\ref{sec:standalone}).

Where mechanisms were inactive or improvements came at a cost, the measurements
are reported as observed and paired with the specific work planned to address
them (Section~\ref{sec:limitations}). Four defects found during the work are
documented in Section~\ref{sec:defects}; all are covered by regression tests, and
the results they invalidated are identified in place.

\begin{thebibliography}{9}
\bibitem{stems2025}
Zhang, X., Wu, J., Zinflou, A., \& Boulet, B. (2025).
\textit{STEMS: Spatial-Temporal Enhanced Multi-Agent Safe Building Energy Management System.}
IEEE Internet of Things Journal. arXiv:2510.14112v2.
\end{thebibliography}

\end{document}
